% ===========================================================================
\chapter{What we are going to build}

Kestrel Island is a small island in the North Atlantic, about four
kilometres across. It has a harbour town, Port Kestrel, a lake, two woods, a
brook, a lighthouse on its northern cape, and roads, lanes, tracks and
footpaths joining them all. It is also entirely imaginary: it exists only
as a file of numbers, \code{kestrel_island.map}.

By the end of this tutorial you will have written a program that turns that
file into a map in a window, and then lets you use it the way you use the
map app on a phone. Click where you want to start and where you want to go.
The program searches for the quickest way, and you can watch the search
spread out across the roads. Then a traveller sets off, a blue dot moving
along the highlighted route, while the program talks them through every
turn: ``In 300 m, turn left onto Mill Lane.''

\mappicture{nav_travel}

Press \code{x}, and the traveller takes a wrong turn on purpose. A moment
later the program notices that they have left the route, works out where
they are, and plans a new route from there, just as a satnav does.

There is a plain-text version too, for any two places:

\laterout{route.nex "beach car park" lighthouse}{route_beach}

\section{How digital maps work}

A map app looks like one program, but it is several quite different ideas
working together. Each of them is simple at heart, and this tutorial builds
every one of them in a few dozen lines of Nex.

\begin{description}[style=nextline, leftmargin=1.2em, itemsep=0.4em]
  \item[Positions on a round Earth.] A GPS gives positions as
    \emph{latitude} and \emph{longitude}: angles, not metres. We need to
    measure distances between them, which means a little spherical
    geometry (Chapter~\ref{ch:latlon}).
  \item[Flattening the Earth.] A screen is flat and the Earth is not.
    A \emph{map projection} turns latitude and longitude into flat
    $x$ and $y$ coordinates in metres. For a small island one simple
    projection is almost perfect (Chapter~\ref{ch:projection}).
  \item[Geometry.] Once things are flat, we need the distance from a point
    to a road, the area of a lake, whether a point is inside a wood, and
    which way a road points (Chapter~\ref{ch:geometry}).
  \item[The map as data.] The map file stores \emph{nodes} (single points)
    and lists of nodes that make roads and areas, the way OpenStreetMap,
    the world's largest open map, does it (Chapter~\ref{ch:mapfile}).
  \item[Drawing.] Turning that data into a picture people can read:
    colours, roads with outlines, labels that don't collide, a scale bar,
    and zooming in and out (Chapters~\ref{ch:drawing}
    and~\ref{ch:labels}).
  \item[Roads as a graph.] For finding a way, a map is a \emph{graph}:
    points joined by one-way or two-way connections, each with a length
    and a travel time (Chapter~\ref{ch:graph}).
  \item[Searching for a route.] Three classic algorithms: breadth-first
    search, Dijkstra's algorithm, and A*, which is what real route
    planners are built on (Chapters~\ref{ch:bfs}--\ref{ch:astar}).
  \item[Navigation.] Turning a route into instructions, following a
    moving traveller from noisy GPS fixes, and planning again when they go
    wrong (Chapters~\ref{ch:directions}--\ref{ch:wrongturn}).
\end{description}

Chapter~\ref{ch:livemap} puts it all together in the live map, and
Chapter~\ref{ch:next} looks at how the real thing differs: maps of whole
countries, tiles, traffic, and the clever tricks that answer a route
question about a continent in a few milliseconds.

\section{How this tutorial works}

We will build the whole project together, from an empty folder, in small
steps. Each step adds a few lines of code to one of the project's files,
and explains what the code does and why. When there is math behind the
code, we work through the math first, in simple steps, and then turn it
into code.

The code comes in boxes like this one. The title of the box says which
file the code belongs to, and where in the file it goes:

\begin{tcolorbox}[codebox, title={\textbf{New file}\enspace\texttt{example.nex}}]
\small Start a new file with this content.
\end{tcolorbox}
\begin{tcolorbox}[codebox, title={\texttt{example.nex}\enspace\textendash\enspace add at the end}]
\small Add this code at the end of the file.
\end{tcolorbox}
\begin{tcolorbox}[codebox, title={\texttt{example.nex}\enspace\textendash\enspace add before the line ...}]
\small Add this code at a particular place in the file, which the title tells you.
\end{tcolorbox}
\begin{tcolorbox}[codebox, title={\texttt{example.nex}\enspace\textendash\enspace continued}]
\small Add this code right after the code from the previous box.
\end{tcolorbox}

Green boxes are \emph{try-it} scripts: small throwaway programs that use
what we just built, so you can see it working. Their output follows in a
dark box, and when they draw a picture, the picture follows too:

\begin{tcolorbox}[trybox, title={\textbf{Try it}\enspace\texttt{try\_example.nex}}]
\small A small program to run now. It is not part of the project.
\end{tcolorbox}
\begin{tcolorbox}[outbox]
{\ttfamily\footnotesize\color{outfg}\$ nex try\_example.nex\\what it prints}
\end{tcolorbox}

Every program, every output and every picture in this tutorial was
produced by running the code exactly as printed here, so if your output
differs, compare your code with the boxes.

\section{What you need}

\begin{itemize}[nosep]
  \item \textbf{Nex}, installed so that the \code{nex} command works in a
        terminal. You should know the basics of the language: \code{let},
        functions and \code{result}, classes, loops, and contracts
        (\code{require}, \code{ensure}, \code{invariant}). We explain the
        less common features as we meet them.
  \item \textbf{The map file}, \code{kestrel_island.map}, from the
        \code{examples/digital_map} folder of the Nex distribution.
  \item \textbf{A display} for the live map in Chapter~\ref{ch:livemap}.
        Everything else, including all the pictures, works without one.
  \item \textbf{Math}: school geometry and trigonometry. We use
        Pythagoras' theorem, sine and cosine, and vectors (arrows), and we
        explain each one when it first comes up.
\end{itemize}

No mapping or graph library is used. Every distance, every projection,
every search and every pixel of the map is computed in Nex, in about 2,500
lines including comments and checks.

\section{The plan}

By the end you will have these files, each of which you will have written
yourself:

\begin{center}
\small
\begin{tabular}{@{}lll@{}}
\toprule
File & What it holds & Chapter \\
\midrule
\code{helpers.nex} & Printing numbers, distances and times neatly & \ref{ch:setup} \\
\code{geo.nex} & Latitude and longitude, flat points, the projection & \ref{ch:latlon}, \ref{ch:projection} \\
\code{geometry.nex} & Lines, segments, polygons and bearings & \ref{ch:geometry} \\
\code{island_map.nex} & Roads, areas and places, read from the map file & \ref{ch:mapfile} \\
\code{map_view.nex} & Drawing the map, at any zoom & \ref{ch:drawing}, \ref{ch:labels} \\
\code{road_graph.nex} & The roads as a graph, for driving or walking & \ref{ch:graph} \\
\code{search.nex} & Breadth-first search, Dijkstra's algorithm and A* & \ref{ch:bfs}--\ref{ch:astar} \\
\code{checks.nex} & Automatic checks of everything & \ref{ch:checks} \\
\code{directions.nex} & Turn-by-turn instructions & \ref{ch:directions} \\
\code{route.nex} & Directions between two places, printed & \ref{ch:directions} \\
\code{navigation.nex} & Following a traveller along a route & \ref{ch:following}, \ref{ch:wrongturn} \\
\code{simulation.nex} & A pretend traveller with a pretend GPS & \ref{ch:wrongturn} \\
\code{navigate.nex} & The live map & \ref{ch:livemap} \\
\bottomrule
\end{tabular}
\end{center}
